\documentclass{ximera}

\input{../../../preamble.tex}


\definecolor{fvdnavy}{HTML}{071D43}
\definecolor{fvdcyan}{HTML}{20BFC6}
\definecolor{fvdcyanlight}{HTML}{EFFCFB}
\definecolor{fvdmagenta}{HTML}{F10D69}
\definecolor{fvdmagentalight}{HTML}{FFF1F7}
\definecolor{fvdtext}{HTML}{163044}





%=================================================
% Pedagogische omgevingen
% PDF: tcolorbox
% HTML: learning-box
%=================================================

\newenvironment{voorkennis}
{%
    \ifdefined\HCode
        \HCode{%
<section class="learning-box learning-box--prior">
<div class="learning-box__header">
<span class="learning-box__icon" aria-hidden="true"></span>
<span class="learning-box__title">Voorkennis</span>
</div>
<div class="learning-box__content">
        }%
        \begin{itemize}
    \else
       \begin{tcolorbox}[
    enhanced,
    breakable,
    colback=fvdcyanlight,
    colframe=fvdcyan,
    coltext=fvdtext,
    boxrule=0.5pt,
    leftrule=4pt,
    arc=4mm,
    outer arc=4mm,
    left=7mm,
    right=6mm,
    top=5mm,
    bottom=5mm,
    before skip=6mm,
    after skip=6mm,
    title=Voorkennis,
    coltitle=fvdcyan,
    fonttitle=\bfseries\Large,
    attach title to upper={\par\smallskip},
    boxed title style={
        empty,
        left=0mm,
        right=0mm,
        top=0mm,
        bottom=0mm
    },
    overlay={
        \draw[
            fvdcyan,
            line width=0.5pt
        ]
        ([xshift=42mm,yshift=-13mm]frame.north west)
        --
        ([xshift=-7mm,yshift=-13mm]frame.north east);
    }
]
        \begin{itemize}
    \fi
}
{%
    \end{itemize}
    \ifdefined\HCode
        \HCode{%
</div>
</section>
        }%
    \else
        \end{tcolorbox}
    \fi
}


\newenvironment{leerdoelen}
{%
    \ifdefined\HCode
        \HCode{%
<section class="learning-box learning-box--goals">
<div class="learning-box__header">
<span class="learning-box__icon" aria-hidden="true"></span>
<span class="learning-box__title">Leerdoelen</span>
</div>
<div class="learning-box__content">
        }%
        \begin{itemize}
    \else
\begin{tcolorbox}[
    enhanced,
    breakable,
    colback=fvdmagentalight,
    colframe=fvdmagenta,
    coltext=fvdtext,
    boxrule=0.5pt,
    leftrule=4pt,
    arc=4mm,
    outer arc=4mm,
    left=7mm,
    right=6mm,
    top=5mm,
    bottom=5mm,
    before skip=6mm,
    after skip=6mm,
    title=Leerdoelen,
    coltitle=fvdmagenta,
    fonttitle=\bfseries\Large,
    attach title to upper={\par\smallskip},
    boxed title style={
        empty,
        left=0mm,
        right=0mm,
        top=0mm,
        bottom=0mm
    },
    overlay={
        \draw[
            fvdmagenta,
            line width=0.5pt
        ]
        ([xshift=42mm,yshift=-13mm]frame.north west)
        --
        ([xshift=-7mm,yshift=-13mm]frame.north east);
    }
]
        \begin{itemize}
    \fi
}
{%
    \end{itemize}
    \ifdefined\HCode
        \HCode{%
</div>
</section>
        }%
    \else
        \end{tcolorbox}
    \fi
}




\title{De afgeleide functie}

\begin{document}

% FVD_CHAPTER_DATA_START
% Automatisch gegenereerd uit index.tex.
% Niet handmatig aanpassen.
\fvdchapterdata{1}{Veranderingen in beeld}{2}
% FVD_CHAPTER_DATA_END




































































































% ============================================================
% ABSTRACT
% ============================================================

\begin{abstract}
In het vorige hoofdstuk bekeken we de afgeleide in één punt:
het afgeleid getal $f'(a)$.

Nu laten we het punt variëren.
Zo ontstaat een nieuwe functie, de afgeleide functie $f'$,
die in elk punt beschrijft hoe snel de oorspronkelijke functie verandert.

We frissen ook op hoe je de eerste en tweede afgeleide van
veeltermfuncties berekent. Die technieken vormen de basis voor het
functieonderzoek in het volgende hoofdstuk en voor de nieuwe
afleidingstechnieken die later in deze module volgen.
\end{abstract}

\maketitle


% ============================================================
% VOORKENNIS EN LEERDOELEN
% ============================================================

\begin{voorkennis}
  \item Je kent het verschil tussen een gemiddelde en een
        ogenblikkelijke verandering.
  \item Je kan het afgeleid getal $f'(a)$ interpreteren.
  \item Je weet dat $f'(a)$ de richtingscoëfficiënt van de raaklijn is.
  \item Je kan werken met machten met gehele positieve exponenten.
  \item Je kan veeltermen vereenvoudigen.
\end{voorkennis}

\begin{leerdoelen}
  \item Je kan het verschil uitleggen tussen een afgeleid getal
        $f'(a)$ en een afgeleide functie $f'(x)$.
  \item Je kan uitleggen wat de functiewaarde $f'(x)$ betekent.
  \item Je kan de afgeleide bepalen van constante functies en
        machtsfuncties $x^n$.
  \item Je kan de constante-factorregel en de som- en verschilregel
        gebruiken.
  \item Je kan veeltermfuncties afleiden.
  \item Je kan de tweede afgeleide $f''$ bepalen.
  \item Je kan de betekenis van $f'$ en $f''$ onderscheiden.
  \item Je kan controleren of een voorgestelde afgeleide plausibel is.
\end{leerdoelen}


%%%%%%%%%%%%%%%%%%%%%%%%%%%%%%%%%%%%%%%%%%%%%%%%%%%%%%%%%%%%%%%%%%%%%%
%% WAAROM LEER JE DIT?
%%%%%%%%%%%%%%%%%%%%%%%%%%%%%%%%%%%%%%%%%%%%%%%%%%%%%%%%%%%%%%%%%%%%%%

\FVDsection{Waarom leer je dit?}

In het vorige hoofdstuk keken we naar de verandering van een functie
in één specifiek punt.

Als bijvoorbeeld

\[
f'(3)=5,
\]

dan weten we hoe $f$ verandert bij $x=3$.

Maar meestal willen we niet alleen weten wat er in één punt gebeurt.
We willen kunnen onderzoeken hoe de verandering zelf varieert
wanneer $x$ verandert.

Daarom stellen we voor elk geschikt punt dezelfde vraag:

\[
\boxed{\text{hoe snel verandert }f\text{ hier?}}
\]

Het antwoord hangt af van $x$ en vormt dus opnieuw een functie.

Die nieuwe functie noemen we de \textbf{afgeleide functie}.

\[
\boxed{
f
\longrightarrow
f'
}
\]

De oorspronkelijke functie $f$ beschrijft de grootheid.

De afgeleide functie $f'$ beschrijft hoe die grootheid verandert.

Later zullen we precies die informatie gebruiken om grafieken
systematisch te onderzoeken.


%%%%%%%%%%%%%%%%%%%%%%%%%%%%%%%%%%%%%%%%%%%%%%%%%%%%%%%%%%%%%%%%%%%%%%
%% VAN AFGELEID GETAL NAAR AFGELEIDE FUNCTIE
%%%%%%%%%%%%%%%%%%%%%%%%%%%%%%%%%%%%%%%%%%%%%%%%%%%%%%%%%%%%%%%%%%%%%%

\FVDsection{Van afgeleid getal naar afgeleide functie}


% --------------------------------------------------------------------
% AFGELEID GETAL
% --------------------------------------------------------------------

\FVDsubsection{Eén punt: het afgeleid getal}

Voor een vaste waarde $a$ is

\[
f'(a)
\]

een \textbf{getal}.

Het beschrijft de ogenblikkelijke verandering van $f$
bij $x=a$.

Meetkundig is het de richtingscoëfficiënt van de raaklijn
aan de grafiek van $f$ in

\[
P(a,f(a)).
\]

Bijvoorbeeld:

\[
f'(2)=4
\]

betekent dat de raaklijn bij $x=2$ richtingscoëfficiënt $4$ heeft.


% --------------------------------------------------------------------
% AFGELEIDE FUNCTIE
% --------------------------------------------------------------------

\FVDsubsection{Alle punten: de afgeleide functie}

We kunnen het afgeleid getal niet alleen bij $x=2$ bepalen,
maar bij elke waarde van $x$ waarvoor de functie afleidbaar is.

Zo ontstaat een nieuwe functie.

\begin{definitie}[Afgeleide functie]

De \textbf{afgeleide functie} van een functie $f$ is de functie

\[
f'
\]

die aan elke waarde $x$ het afgeleid getal van $f$ in dat punt
toekent.

Dus:

\[
\boxed{
x\longmapsto f'(x)
}
\]

waarbij $f'(x)$ de ogenblikkelijke verandering van $f$
bij $x$ beschrijft.

\end{definitie}


Het onderscheid is belangrijk:

\[
\boxed{f'(a)}
\qquad\text{is een getal,}
\]

terwijl

\[
\boxed{f'}
\qquad\text{een functie is.}
\]


\begin{voorbeeld}

Beschouw

\[
f(x)=x^2.
\]

De afgeleide functie is

\[
f'(x)=2x.
\]

Daaruit kunnen we verschillende afgeleide getallen bepalen:

\[
f'(1)=2,
\qquad
f'(3)=6,
\qquad
f'(-2)=-4.
\]

De functie

\[
f'(x)=2x
\]

geeft dus voor elk punt onmiddellijk de richtingscoëfficiënt
van de raaklijn aan de grafiek van $f$.

\end{voorbeeld}


\begin{oefening}[Afgeleid getal of afgeleide functie?]{\oefB \ESS}

Gegeven is

\[
f'(x)=3x^2-4.
\]

\begin{question}

Is $f'$ een getal of een functie?

\begin{multipleChoice}
  \choice{Een getal}
  \choice[correct]{Een functie}
\end{multipleChoice}

\end{question}

\begin{question}

Bereken

\[
f'(2)=\answer{8}.
\]

\end{question}

\begin{question}

Wat stelt $f'(2)$ meetkundig voor?

\begin{multipleChoice}
  \choice{De $y$-coördinaat van een punt op de grafiek van $f$.}
  \choice[correct]{De richtingscoëfficiënt van de raaklijn aan de grafiek van $f$ bij $x=2$.}
  \choice{De gemiddelde verandering van $f$ over $[0,2]$.}
\end{multipleChoice}

\end{question}

\end{oefening}


%%%%%%%%%%%%%%%%%%%%%%%%%%%%%%%%%%%%%%%%%%%%%%%%%%%%%%%%%%%%%%%%%%%%%%
%% BASISAFGELEIDEN
%%%%%%%%%%%%%%%%%%%%%%%%%%%%%%%%%%%%%%%%%%%%%%%%%%%%%%%%%%%%%%%%%%%%%%

\FVDsection{Basisregels voor veeltermfuncties}

We frissen nu de afleidingsregels op die nodig zijn om
veeltermfuncties af te leiden.



% --------------------------------------------------------------------
% CONSTANTE FUNCTIE
% --------------------------------------------------------------------

\FVDsubsection{De afgeleide van een constante functie}

Beschouw een constante functie

\[
f(x)=c.
\]

De grafiek is een horizontale rechte.

De richtingscoëfficiënt is overal nul.

Daarom geldt:

\begin{eigenschap}[Constante functie]

Voor elke constante $c$:

\[
\boxed{
(c)'=0
}
\]

of equivalent:

\[
f(x)=c
\quad\Longrightarrow\quad
f'(x)=0.
\]

\end{eigenschap}


\begin{voorbeeld}

Als

\[
f(x)=7,
\]

dan is

\[
f'(x)=0.
\]

Ook voor

\[
g(x)=-13
\]

geldt

\[
g'(x)=0.
\]

\end{voorbeeld}


% --------------------------------------------------------------------
% MACHTSREGEL
% --------------------------------------------------------------------

\FVDsubsection{De machtsregel}

Een van de belangrijkste basisregels is de machtsregel.

\begin{eigenschap}[Machtsregel]

Voor een positieve gehele exponent $n$ geldt:

\[
\boxed{
(x^n)'=n x^{n-1}
}
\]

\end{eigenschap}


Je kunt de regel onthouden als:

\[
\boxed{
\text{exponent naar voren}
\quad\text{en}\quad
\text{exponent één verminderen}.
}
\]


\begin{voorbeeld}

\[
f(x)=x^5
\]

geeft

\[
f'(x)=5x^4.
\]

Voor

\[
g(x)=x^3
\]

vinden we

\[
g'(x)=3x^2.
\]

En voor

\[
h(x)=x
\]

kunnen we schrijven

\[
h(x)=x^1,
\]

zodat

\[
h'(x)=1x^0=1.
\]

\end{voorbeeld}


\begin{opmerking}

Bij

\[
x^1
\]

wordt de exponent na afleiden nul.

Omdat

\[
x^0=1,
\]

geldt dus

\[
(x)'=1.
\]

\end{opmerking}


% --------------------------------------------------------------------
% CONSTANTE FACTOR
% --------------------------------------------------------------------

\FVDsubsection{Een constante factor}

Een constante factor blijft bij het afleiden behouden.

\begin{eigenschap}[Constante-factorregel]

Als $c$ een constante is, dan:

\[
\boxed{
(c\,f(x))'=c\,f'(x)
}
\]

\end{eigenschap}


In het bijzonder:

\[
(ax^n)'=a n x^{n-1}.
\]


\begin{voorbeeld}

Voor

\[
f(x)=4x^5
\]

geldt

\[
f'(x)=4\cdot5x^4=20x^4.
\]

Voor

\[
g(x)=-3x^2
\]

geldt

\[
g'(x)=-3\cdot2x=-6x.
\]

\end{voorbeeld}


% --------------------------------------------------------------------
% SOM EN VERSCHIL
% --------------------------------------------------------------------

\FVDsubsection{Som- en verschilregel}

Bij een som of verschil mogen we term per term afleiden.

\begin{eigenschap}[Som- en verschilregel]

Voor afleidbare functies $f$ en $g$ geldt:

\[
\boxed{
(f+g)'=f'+g'
}
\]

en

\[
\boxed{
(f-g)'=f'-g'.
}
\]

\end{eigenschap}


Daardoor kunnen we een veelterm eenvoudig term per term afleiden.


\begin{voorbeeld}

Gegeven:

\[
f(x)=3x^4-5x^2+7x-9.
\]

We leiden elke term afzonderlijk af:

\[
f'(x)
=
12x^3-10x+7.
\]

De constante term $-9$ verdwijnt omdat

\[
(-9)'=0.
\]

\end{voorbeeld}


%%%%%%%%%%%%%%%%%%%%%%%%%%%%%%%%%%%%%%%%%%%%%%%%%%%%%%%%%%%%%%%%%%%%%%
%% VEELTERMFUNCTIES
%%%%%%%%%%%%%%%%%%%%%%%%%%%%%%%%%%%%%%%%%%%%%%%%%%%%%%%%%%%%%%%%%%%%%%

\FVDsection{Veeltermfuncties afleiden}


% --------------------------------------------------------------------
% ALGEMENE VORM
% --------------------------------------------------------------------

\FVDsubsection{Term per term}

Voor een veeltermfunctie

\[
f(x)
=
a_nx^n+a_{n-1}x^{n-1}+\cdots+a_1x+a_0
\]

kunnen we elke term afzonderlijk afleiden.

We krijgen:

\[
\boxed{
f'(x)
=
na_nx^{n-1}
+(n-1)a_{n-1}x^{n-2}
+\cdots
+a_1.
}
\]


\begin{voorbeeld}

Bepaal de afgeleide van

\[
f(x)
=
-2x^5+3x^4-7x^2+6.
\]

\begin{oplossing}[toon]

We leiden term per term af:

\[
(-2x^5)'=-10x^4,
\]

\[
(3x^4)'=12x^3,
\]

\[
(-7x^2)'=-14x,
\]

en

\[
(6)'=0.
\]

Dus:

\[
\boxed{
f'(x)
=
-10x^4+12x^3-14x.
}
\]

\end{oplossing}

\end{voorbeeld}


\begin{oefening}[Veeltermfuncties afleiden]{\oefB \ESS}

Bepaal telkens de afgeleide functie.

\begin{question}
\(
f(x)=5x^4-3x^2+8
\) \ \ \ \ \ \ \ \ 
\(
f'(x)=\answer{20x^3-6x}.
\)

\end{question}

\begin{question}
\(
g(x)=-2x^3+7x-4
\) \ \ \ \ \ \ \ \ 
\(
g'(x)=\answer{-6x^2+7}.
\)

\end{question}

\begin{question}
\(
h(x)=\frac{1}{2}x^6-4x^3+x
\) \ \ \ \ \ \ \ \ 
\(
h'(x)=\answer{3x^5-12x^2+1}.
\)

\end{question}

\end{oefening}


% --------------------------------------------------------------------
% EERST VEREENVOUDIGEN
% --------------------------------------------------------------------

\FVDsubsection{Vereenvoudig wanneer dat nuttig is}

Een functie staat niet altijd onmiddellijk in de handigste vorm.

Bij veeltermfuncties kan het nuttig zijn om eerst haakjes weg te werken
of gelijksoortige termen samen te nemen.

\begin{voorbeeld}

Gegeven:

\[
f(x)=2x(x^2-3x+4).
\]

Omdat dit een product van veeltermen is, kunnen we eerst uitwerken:

\[
f(x)=2x^3-6x^2+8x.
\]

Nu leiden we term per term af:

\[
\boxed{
f'(x)=6x^2-12x+8.
}
\]

\end{voorbeeld}


\begin{oefening}[Eerst vereenvoudigen]{\oefU \ESS}

Gegeven:

\[
f(x)=(x+2)(x^2-3x+1).
\]

Werk eerst uit en bepaal daarna $f'(x)$.

\begin{oplossing}

Eerst:

\[
f(x)
=
x^3-x^2-5x+2.
\]

Daarom:

\[
\boxed{
f'(x)=3x^2-2x-5.
}
\]

\end{oplossing}

\end{oefening}


%%%%%%%%%%%%%%%%%%%%%%%%%%%%%%%%%%%%%%%%%%%%%%%%%%%%%%%%%%%%%%%%%%%%%%
%% EEN AFGELEIDE CONTROLEREN
%%%%%%%%%%%%%%%%%%%%%%%%%%%%%%%%%%%%%%%%%%%%%%%%%%%%%%%%%%%%%%%%%%%%%%

\FVDsection{Controleer je afgeleide}


% --------------------------------------------------------------------
% GRAAD
% --------------------------------------------------------------------

\FVDsubsection{De graad daalt met één}

Bij een niet-constante veeltermfunctie daalt de graad na afleiden
met één.

Bijvoorbeeld:

\[
f(x)=4x^5-2x^3+x
\]

heeft graad $5$.

De afgeleide

\[
f'(x)=20x^4-6x^2+1
\]

heeft graad $4$.

Dit geeft een snelle controle.


% --------------------------------------------------------------------
% CONSTANTEN
% --------------------------------------------------------------------

\FVDsubsection{Constante termen verdwijnen}

Een constante term heeft geen invloed op de verandering van een functie.

Daarom verdwijnt hij bij het afleiden.

Als

\[
f(x)=3x^2+5
\]

en

\[
g(x)=3x^2-100,
\]

dan geldt voor beide:

\[
f'(x)=g'(x)=6x.
\]

De grafieken liggen verticaal verschoven ten opzichte van elkaar,
maar hebben bij dezelfde $x$ dezelfde helling.


\begin{oefening}[Welke afgeleide kan kloppen?]{\oefU \ESS}

Gegeven is

\[
f(x)=2x^4-5x^2+7.
\]

Welke uitdrukking kan de afgeleide zijn?

\begin{multipleChoice}

\choice{
$8x^4-10x$
}

\choice[correct]{
$8x^3-10x$
}

\choice{
$8x^3-10x+7$
}

\choice{
$2x^3-5x$
}

\end{multipleChoice}

Leg je keuze uit aan de hand van minstens twee afleidingsregels.

\begin{oplossing}

  We leiden de functie term per term af:
  
  \[
  f(x)=2x^4-5x^2+7.
  \]
  
  Met de \textbf{machtsregel}
  
  \[
  (x^n)'=nx^{n-1}
  \]
  
  en de \textbf{constante-factorregel}
  
  \[
  (c\cdot f(x))'=c\cdot f'(x)
  \]
  
  vinden we
  
  \[
  (2x^4)'=2\cdot 4x^3=8x^3
  \]
  
  en
  
  \[
  (-5x^2)'=-5\cdot 2x=-10x.
  \]
  
  De afgeleide van een constante is nul, dus
  
  \[
  (7)'=0.
  \]
  
  Daarom is
  
  \[
  \boxed{f'(x)=8x^3-10x}.
  \]
  
  De tweede uitdrukking is dus de enige die kan kloppen.
  
  \end{oplossing}

\end{oefening}


%%%%%%%%%%%%%%%%%%%%%%%%%%%%%%%%%%%%%%%%%%%%%%%%%%%%%%%%%%%%%%%%%%%%%%
%% EERSTE EN TWEEDE AFGELEIDE
%%%%%%%%%%%%%%%%%%%%%%%%%%%%%%%%%%%%%%%%%%%%%%%%%%%%%%%%%%%%%%%%%%%%%%

\FVDsection{De eerste en tweede afgeleide}


% --------------------------------------------------------------------
% TWEEDE AFGELEIDE
% --------------------------------------------------------------------

\FVDsubsection{De afgeleide opnieuw afleiden}

Omdat $f'$ zelf een functie is, kunnen we — wanneer dit mogelijk is —
ook $f'$ opnieuw afleiden.

Zo ontstaat de tweede afgeleide.

\begin{definitie}[Tweede afgeleide]

De afgeleide van de afgeleide functie noemen we
de \textbf{tweede afgeleide} van $f$.

We noteren:

\[
\boxed{
f''(x)=(f'(x))'.
}
\]

\end{definitie}


\begin{voorbeeld}

Gegeven:

\[
f(x)=2x^4-3x^3+5x.
\]

De eerste afgeleide is

\[
f'(x)=8x^3-9x^2+5.
\]

We leiden opnieuw af:

\[
f''(x)=24x^2-18x.
\]

Dus:

\[
\boxed{
f'(x)=8x^3-9x^2+5
}
\]

en

\[
\boxed{
f''(x)=24x^2-18x.
}
\]

\end{voorbeeld}


\begin{oefening}[Eerste en tweede afgeleide]{\oefB \ESS}

Gegeven:

\[
f(x)=x^5-4x^3+6x^2-2.
\]

Bepaal $f'$ en $f''$.

\begin{itemize}
  \item \( f'(x)=\answer{5x^4-12x^2+12x} \)
  \item \(f''(x)=\answer{20x^3-24x+12} \)
\end{itemize}

\end{oefening}


\begin{voorbeeld}[Beweging]

Stel dat

\[
s(t)
\]

de positie van een bewegend voorwerp beschrijft in meter,
met $t$ in seconde.

Dan is

\[
s'(t)
\]

de ogenblikkelijke snelheid in

\[
\mathrm{m/s},
\]

en

\[
s''(t)
\]

de ogenblikkelijke versnelling in

\[
\mathrm{m/s^2}.
\]

Hier zien we zeer concreet het verschil tussen een functie,
haar eerste afgeleide en haar tweede afgeleide.

\end{voorbeeld}


\begin{oefening}[Beweging]{\oefU \ESS}

De positie van een voorwerp wordt gegeven door

\[
s(t)=t^3-6t^2+9t
\]

met $s$ in meter en $t$ in seconde.

\begin{question}

Bepaal de snelheidsfunctie.

\[
v(t)=s'(t)=\answer{3t^2-12t+9}.
\]

\end{question}

\begin{question}

Bepaal de versnellingsfunctie.

\[
a(t)=s''(t)=\answer{6t-12}.
\]

\end{question}

\begin{question}

Bereken de snelheid bij $t=2$.

\[
v(2)=\answer{-3}.
\]

\end{question}

\begin{question}

Welke eenheid hoort bij $v(2)$?

\begin{multipleChoice}
  \choice{$\mathrm{m}$}
  \choice[correct]{$\mathrm{m/s}$}
  \choice{$\mathrm{m/s^2}$}
\end{multipleChoice}

\end{question}

\begin{question}

Bereken de versnelling bij $t=2$.

\[
a(2)=\answer{0}.
\]

\end{question}

\end{oefening}







%%%%%%%%%%%%%%%%%%%%%%%%%%%%%%%%%%%%%%%%%%%%%%%%%%%%%%%%%%%%%%%%%%%%%%
%% SAMENVATTING
%%%%%%%%%%%%%%%%%%%%%%%%%%%%%%%%%%%%%%%%%%%%%%%%%%%%%%%%%%%%%%%%%%%%%%

\FVDsection{Samengevat}

Het afgeleid getal

\[
f'(a)
\]

beschrijft de ogenblikkelijke verandering in één punt.

Wanneer we dit voor elke geschikte waarde van $x$ doen,
krijgen we de afgeleide functie

\[
\boxed{
f'(x).
}
\]

Voor veeltermfuncties gebruiken we vooral:

\[
\boxed{
(c)'=0
}
\]

\[
\boxed{
(x^n)'=nx^{n-1}
}
\]

\[
\boxed{
(c f)'=c f'
}
\]

\[
\boxed{
(f\pm g)'=f'\pm g'.
}
\]

De afgeleide functie kan opnieuw worden afgeleid:

\[
\boxed{
f''=(f')'.
}
\]

\begin{oefening}[Van functie naar verandering]{\oefU \ESS}

  Gegeven is
  
  \[
  f(x)=x^3-3x^2-9x+5.
  \]
  
  \begin{question}
  
  Bepaal de eerste afgeleide.
  
  \[
  f'(x)=\answer{3x^2-6x-9}.
  \]
  
  \end{question}
  
  \begin{question}
  
  Bepaal de tweede afgeleide.
  
  \[
  f''(x)=\answer{6x-6}.
  \]
  
  \end{question}
  
  \begin{question}
  
  Bereken
  
  \[
  f'(0)=\answer{-9}.
  \]
  
  Wat betekent dit meetkundig?
  
  \begin{multipleChoice}
    \choice{De grafiek gaat door $(0,-9)$.}
    \choice[correct]{De raaklijn aan de grafiek bij $x=0$ heeft richtingscoëfficiënt $-9$.}
    \choice{De functie heeft noodzakelijk een minimum bij $x=0$.}
  \end{multipleChoice}
  
  \end{question}
  
  \begin{question}
  
  Een leerling beweert:
  
  \begin{center}
  ``Omdat $f'(0)<0$, is $f$ overal dalend.''
  \end{center}
  
  Is dat besluit correct?
  
  \begin{multipleChoice}
  
  \choice{
  Ja, want een negatieve afgeleide betekent altijd dat een functie
  over haar volledige domein daalt.
  }
  
  \choice[correct]{
  Nee. $f'(0)<0$ geeft alleen informatie over de verandering
  bij $x=0$. Om te weten waar $f$ stijgt of daalt, moeten we
  het teken van $f'(x)$ over intervallen onderzoeken.
  }
  
  \end{multipleChoice}
  
  \end{question}
  
  \end{oefening}
  
  
  \begin{oefening}[Redeneren zonder volledig uit te rekenen]{\oefG}
  
  Een veeltermfunctie $f$ heeft graad $6$ en een niet-nul
  coëfficiënt bij $x^6$.
  
  Beantwoord zonder een concreet functievoorschrift te kiezen.
  
  \begin{question}
  
  Welke graad heeft $f'$?
  
  \[
  \answer{5}
  \]
  
  \end{question}
  
  \begin{question}
  
  Welke graad heeft $f''$?
  
  \[
  \answer{4}
  \]
  
  \end{question}
  
  \begin{question}
  
  Kan $f'$ een constante functie zijn?
  
  \begin{multipleChoice}
    \choice{Ja}
    \choice[correct]{Nee}
  \end{multipleChoice}

  \begin{oplossing}
    De afgeleide $f'$ kan dus geen constante functie zijn. Een constante
    functie heeft graad $0$, terwijl $f'$ hier noodzakelijk graad $5$ heeft.
    
    \[
    \boxed{\text{$f'$ kan geen constante functie zijn.}}
    \]
  \end{oplossing}
  
  \end{question}
  
  Leg uit waarom je antwoorden noodzakelijk volgen uit de machtsregel.

  \begin{oplossing}

Algemeen verlaagt het afleiden van een niet-constante veelterm de graad
telkens met één. Dit volgt rechtstreeks uit de machtsregel.


  \end{oplossing}
  
  \end{oefening}
  




\end{document}